\documentclass[10pt,a4paper]{amsart}
\usepackage[margin=19mm]{geometry}
\usepackage{amsmath,amssymb,mathtools}
\usepackage[scr=rsfs,cal=euler]{mathalfa}
\usepackage{xfrac}
\usepackage[unicode,hidelinks,bookmarks=false]{hyperref}
\newcommand{\ie}{i.e.\ }
\newcommand{\cf}{cf.\ }
\newcommand{\resp}{respectively\ }
\newcommand{\Subsection}{\subsection}
\providecommand{\nobreakdash}{\nobreak\hbox{-}}
\setlength{\emergencystretch}{2em}
\multlinegap0pt
\usepackage{bm}
\usepackage{bbm}
\usepackage{dsfont}
\usepackage{xspace}
\usepackage{cancel}
\usepackage{mathtools}
\usepackage{amsthm}
\makeatletter
\@ifundefined{thm}{\newtheorem{thm}{Theorem}}{}
\makeatother
\title[The Cayley-Moser problem with Poissonian arrival of offers]{The Cayley-Moser problem with Poissonian arrival of offers}
\author{Guy Katriel}
\date{Source: arXiv:2511.02763v1 (CC BY 4.0)}
\begin{document}
\maketitle

\noindent\emph{Source and licence.} The Cayley-Moser problem with Poissonian arrival of offers. Version arXiv:2511.02763v1 (CC BY 4.0), distributed under the Creative Commons Attribution 4.0 International licence, as recorded in the arXiv licence metadata for this exact version and shown on the abstract page https://arxiv.org/abs/2511.02763v1. Full rights notice: licenses/arxiv-2511.02763-CC-BY-4.0.txt.\par\smallskip

\noindent\textit{Editorial scope.} Counted unit: the Thm whose source label is \texttt{char} in the article (source0.tex lines 535--563, source label \texttt{char}), delivered together with the article's own complete original proof of it (source0.tex lines 565--698, one \texttt{proof} environment of 134 lines). No mathematics is written, repaired or re-derived: statement and proof are the article's own text line for line. Required context retained uncounted: fm (lines 215--216); def\_M (lines 218--218); fm0 (lines 226--227); mM (lines 232--232); defphi (lines 315--316); ef (lines 329--329); def\_Psi (lines 331--331). Every definition, display and label of those lines is retained unchanged. Omitted from this package and not redistributed: 1--214, 217--217, 219--225, 228--231, 233--314, 317--328, 330--330, 332--534, 564--564, 699--2303. Nothing inside a retained excerpt is omitted or altered: every retained body is delivered exactly as the article's lines are written, so no omission receipt of any kind accompanies this package. Citations: the retained excerpts carry no citation command at all, so no bibliography entry of the article is transcribed and no reference is redistributed. Reference adaptation: none was needed and none was made; every reference inside the delivered unit resolves inside it, so no unresolved reference or citation is produced. Printed slips kept literal: no printed word, symbol, hypothesis or number of the counted statement or of its proof was altered. Layout and class: the article's own class is \texttt{article} with packages including amsmath, amsthm, amssymb, graphicx, fullpage, listings, xcolor, hyperref; the wrapper is the lane's pinned \texttt{amsart} editorial wrapper at 10pt on A4, which restates the theorem-like environments the retained text needs and adds the article's own macro definitions that the retained text needs (none), with extra wrapper packages (bm, bbm, dsfont, xspace, cancel, mathtools). The delivered numbering is the unit's own. Assets: no retained span contains a figure environment, a graphics-inclusion command, a PDF-page inclusion, a table or a TikZ picture; no image file is copied into this package, the package declares no asset dependency and no image file is redistributed with it. Licence: version arXiv:2511.02763v1 of this article is distributed under the Creative Commons Attribution 4.0 International licence, as recorded in the arXiv licence metadata for this exact version and shown on the abstract page https://arxiv.org/abs/2511.02763v1. A full rights notice is delivered at licenses/arxiv-2511.02763-CC-BY-4.0.txt. Characters: every retained excerpt is pure ASCII, so the delivered engine, XeLaTeX with its default Latin Modern text font, reports no missing character.
	\begin{equation}\label{fm}
		E[X_+]=\int_{0}^\infty xdF(x)<\infty.\end{equation}

	\begin{equation}\label{def_M}M=\sup\; \{ x\;|\; F(x)<1\},\end{equation}

	\begin{equation}\label{fm0} E[|X_0|]=\int_{-\infty}^\infty |x|dF_0(x)<\infty.
	\end{equation}

	\begin{equation}\label{mM}0\leq \mu_0<M,\end{equation}

\begin{equation} \label{defphi}\phi(x)=\int_x^\infty [1-F(u)]du,
\end{equation}

	\begin{equation}\label{ef}\mu(t)=\Psi^{-1}(\lambda t),\end{equation}

	\begin{equation}\label{def_Psi}\Psi(x)\doteq\int_{\mu_0}^x \frac{du}{\phi(u)},\;\;\;\; x\in [\mu_0,M).\end{equation}

	\begin{thm}\label{char}
		(1) Let $F,F_0$ be a pair of distributions satisfying \eqref{fm},\eqref{fm0},\eqref{mM}, and $\lambda>0$.
		Then the corresponding optimal policy $\mu: [0,\infty)\rightarrow [0,\infty)$ has the following properties:
		\begin{itemize}
			\item[(i)] $\mu(0)=\mu_0=\int_{-\infty}^\infty xdF_0(x)$ and  $\lim_{t\rightarrow +\infty}\mu(t)=M$, where $M$ is defined by
			\eqref{def_M}.
			
			\item[(ii)] $\mu$ is continuously differentiable.
			
			\item[(iii)] $\mu'$ is absolutely continuous and satisfies $\mu'(t)>0$ for all $t\geq 0$. In particular $\mu$ is strictly monotone increasing.
			
			\item[(iv)] $\mu(t)$ is strictly concave.
			
			\item[(v)] The function \begin{equation}\label{def_h}h(t)=-\frac{\mu''(t)}{\mu'(t)}
			\end{equation}
			(which, by (iii),(iv) is defined and positive {\it{a.e.}} on $[0,\infty)$) is a monotone decreasing function.
			
			\item[(vi)] The function $h$ satisfies
			\begin{equation}\label{h_cond}h(0+)\doteq\lim_{t\rightarrow 0+}h(t)\leq \lambda.
			\end{equation}
		\end{itemize}
		
		(2) Conversely, let $\mu:[0,\infty)\rightarrow [0,\infty)$ be any function
		satisfying (ii),(iii),(iv),(v) and let $\lambda$ satisfy \eqref{h_cond}.  Then, defining
		$M=\lim_{t\rightarrow \infty}\mu(t),$ there 
		exists a distribution $F$ satisfying 
		\eqref{fm}, such that for any distribution $F_0$ satisfying \eqref{fm0} with $E[X_0]=\mu(0)$, we have that 
		$\mu(t)$ is the optimal policy corresponding to $F,F_0$, and $\lambda$.
	\end{thm}

	\begin{proof}
		By its definition \eqref{defphi}, the function $\phi$ is positive on $(-\infty,M)$. In addition
		$\phi$ is absolutely continuous  $[0,M)$, with
		$$\phi'(x)=F(x)-1<0,\;\;\;{\mbox{for {\it{a.e.}} }}\;x\in [0,M),$$
		and in particular is strictly decreasing on 
		$(-\infty,M)$.
		We also have (both when $M<\infty$ and when $M=+\infty$)
		\begin{equation}\label{pgz}
			\lim_{x\rightarrow M-}\phi(x)=0.
		\end{equation}
		
		By our assumption \eqref{mM} and the positivity of $\phi$, the function
		$\Psi$ given by \eqref{def_Psi} is well-defined on $[\mu_0,M)$.
		Moreover $\Psi$ is continuously differentiable on $[\mu_0,M)$, with
		\begin{equation}\label{dpsip}\Psi'(x)=\frac{1}{\phi(x)}>0,\end{equation}
		and in particular it is strictly monotone
		increasing on $[\mu_0,M)$. Since $\phi(x)$ is absolutely continuous and positive 
		$\Psi'(x)=\frac{1}{\phi(x)}$ is also absolutely continuous on any closed subinterval of $[\mu_0,M)$.
		If $M<\infty$, then $\phi(M)=0$ and
		\begin{align*}\phi(u)&=\int_{u}^M  [1-F(x)]dx\leq M-u,\;\;\; u<M\\
			&\Rightarrow\;\;\Psi(x)=\int_{\mu_0}^x \frac{du}{\phi(u)}\geq \int_{\mu_0}^x
			\frac{du}{M-u}=\ln\left( \frac{M-\mu_0}{M-x}\right),
		\end{align*}
		which implies 
		\begin{equation}\label{psi_inf}\lim_{x\rightarrow M-}\Psi(x)=+\infty.\end{equation}
		\eqref{psi_inf} also holds if $M=+\infty$,  
		since, in this case, by \eqref{pgz} the integrand $1/\phi(u)$ in \eqref{def_Psi} 
		goes to $+\infty$ as $u\rightarrow\infty$.
		Since $\phi(x)$ is strictly decreasing, 
		$\Psi'(x)=\frac{1}{\phi(x)}$ is strictly increasing, hence $\Psi(x)$ is strictly convex.
		
		Since $\Psi$  satisfies $\Psi(\mu_0)=0$, and by \eqref{dpsip}, and \eqref{psi_inf},
		the function 
		$\Psi^{-1}$ is well defined and strictly increasing on $[0,\infty)$, with
		\begin{equation}\label{psip}\Psi^{-1}(0)=\mu_0,\;\;\; \lim_{v\rightarrow +\infty} \Psi^{-1}(v)=M,\end{equation}
		and is also continuously differentiable.
		Moreover, since 
		$$(\Psi^{-1})'(z)=\frac{1}{\Psi'(\Psi^{-1}(z))}=\phi(\Psi^{-1}(z)),\;\;\;z\in [0,\infty),$$
		we have that $(\Psi^{-1})'(z)$ is positive and is absolutely continuous on every closed sub-interval of $[0,\infty)$, as a composition of the absolutely continuous function $\phi$ and the continuously  differentiable function $\Psi^{-1}$.
		Since $\Psi$ is strictly convex 
		$\Psi^{-1}$ is strictly concave.
		
		Thus, by \eqref{ef} and \eqref{psip}, $\mu(t)$ is defined on $[0,\infty)$, with
		\begin{equation}\label{ll}\mu(0)=\mu_0,\;\;\;\lim_{t\rightarrow +\infty}\mu(t)=M,\end{equation}
		so we have (i).
		Since $\Psi^{-1}$ is continuously differentiable and $(\Psi^{-1})'$ is positive and absolutely continuous, \eqref{ef} implies 
		(ii),(iii). Since $\Psi^{-1}$ is strictly concave, 
		\eqref{ef}, implies the same for $\mu(t)$, so we have (iv). Finally, note that we have,
		for {\it{a.e.}} $t\geq 0$,
		\begin{equation}\label{prop}
			\mu'(t)=\frac{\lambda}{\Psi'(\Psi^{-1}(\lambda t))}=\lambda\phi(\mu(t))\;\;\Rightarrow\;\; \mu''(t)=\lambda \phi'(\mu(t))\mu'(t)\;\;\Rightarrow\;\; 
			h(t)=-\frac{\mu''(t)}{\mu'(t)}=\lambda[1-F(\mu(t))],\end{equation}
		and since $F$, $\mu$ are increasing functions, we conclude that $h$
		is monotone decreasing. Also, by \eqref{prop}, we have
		$h(0+)=\lambda (1-F(\mu_0+))\leq \lambda$, so that \eqref{h_cond} holds.
		
		(2) To prove the converse direction, assume 
		$\mu: [0,\infty)\rightarrow [0,\infty)$  is a function satisfying  (ii),(iii),(iv), and define
		\begin{equation}\label{dd}
			\mu_0=\mu(0),\;\;M=\lim_{t\rightarrow +\infty}\mu(t).
		\end{equation}
		We also assume $\lambda$ satisfies \eqref{h_cond}. 
		We will construct 
		distribution a $F(x)$ satisfying 
		\eqref{fm}, such that, for
		an arbitrary 
		distribution $F_0(x)$  with $\int_{-\infty}^\infty xdF_0(x)dx=\mu_0$,
		$\mu(t)$ is the optimal policy corresponding to $F$ and $F_0$.
		
		By \eqref{dd} and (iii), we have that $\mu_0<M$ and that the function 
		$\mu^{-1}:[\mu_0,M)\rightarrow [0,\infty)$ exists, is strictly monotone increasing, and satisfies
		\begin{equation}\label{mip}\mu^{-1}(\mu_0)=0,\;\;\lim_{x\rightarrow M-}\mu^{-1}(x)=+\infty.\end{equation}
		
		By (ii),(iii),(v) we have that the function $h(t)$ defined by \eqref{def_h} is an {\it{a.e.}} defined, non-negative and monotone decreasing function, hence  the 
		limit 
		$$h(+\infty)=\lim_{t\rightarrow +\infty}h(t)$$
		exists and satisfies $h(+\infty)\geq 0$. 
		We claim that
		\begin{equation}\label{claim}{\mbox{if }}M=+\infty\;\;{\mbox{then}}\;\;h(+\infty)=0.
		\end{equation}
		Indeed, since
		$$\ln\left(\frac{\mu'(t)}{\mu'(0)}\right)=-\int_0^t h(s)ds\leq -\int_0^t h(+\infty)dt=-h(+\infty)t\;\;\Rightarrow\;\; \mu'(t)\leq \mu'(0)e^{-h(+\infty)t},$$
		if $h(+\infty)>0$ we obtain
		$$\mu(t)\leq \mu(0)+\mu'(0)\int_0^t e^{-h(+\infty)s}ds\leq \mu(0)+\mu'(0)\int_0^{\infty} e^{-h(+\infty)s}ds <\infty,$$
		so $M<\infty$, and we have proved the claim \eqref{claim}.
		
		We now define $F(x)$ by
		\begin{equation}\label{def_FF}F(x)=\begin{cases}
				0 & x<\mu_0\\
				1-\frac{1}{\lambda}h(\mu^{-1}(x)) & \mu_0\leq x<M\\
				1 & x\geq M
		\end{cases}\end{equation}
		(if $M=+\infty$ then the last term for $x\geq M$ is superfluous). To show that $F(x)$ is 
		a cumulative distribution function, we have to show it is monotone increasing  
		with $F(-\infty)=0,F(+\infty)=1$.
		Since $h$ is decreasing and $\mu^{-1}$ is increasing, 
		$F$ is a monotone increasing function on
		$[\mu_0,M)$, and since 
		$\lambda$ satisfies \eqref{h_cond}, we have
		\begin{equation*}\label{gi0}F(\mu_0+)=1-\frac{1}{\lambda} h(0+)\geq 0= F(\mu_0-).
		\end{equation*}
		Using \eqref{mip} and the fact that $h(+\infty)\leq 0$, we have
		\begin{equation}\label{gi}
			\lim_{x\rightarrow M-} F(x)=1-\frac{1}{\lambda}\cdot
			\lim_{x\rightarrow M-}h(\mu^{-1}(x))=1-\frac{1}{\lambda}\cdot
			h(+\infty)\leq 1.
		\end{equation}
		Therefore $F$ is monotone increasing on 
		$(-\infty,\infty)$. 
		Moreover, if $M<\infty$ then we have
		$F(+\infty)=1$ by \eqref{def_FF}, and if $M=+\infty$ then, by 
		\eqref{claim}, we have $h(+\infty)=0$, hence \eqref{gi} 
		becomes an equality, which again gives $F(+\infty)=1$. 
		We have therefore shown that $F(x)$ 
		is a cumulative probability distribution.
		
		It remains to show that $\mu(t)$ is 
		the optimal policy corresponding to the distribution $F$ we defined. Let us compute 
		the optimal policy $\tilde{\mu}(t)$  corresponding to $F,F_0$ (recall $F_0$ is an arbitrary distribution with expectation $\mu_0=\mu(0)$), and show that $\tilde{\mu}(t)=\mu(t)$.
		
		Using
		\eqref{def_FF}, and making the change of variable $u=\mu(t),\;\;du=\mu'(t)dt$, we have, for 
		$x\in [\mu_0,M)$,
		\begin{align*}\phi(x)&=\int_x^\infty (1-F(u))du
			=\frac{1}{\lambda}\int_x^M h(\mu^{-1}(u))du=\frac{1}{\lambda}\int_{\mu^{-1}(x)}^{\infty} h(t)\mu'(t)dt\\&=\frac{1}{\lambda}\int_{\mu^{-1}(x)}^{\infty} \mu''(t)dt=\frac{1}{\lambda}\mu'(\mu^{-1}(x)),
		\end{align*}
		$$\Psi(x)=\int_{\mu_0}^x \frac{du}{\phi(u)}=\lambda\int_{\mu_0}^x \frac{du}{\mu'(\mu^{-1}(u))}=\lambda\int_{\mu_0}^x [\mu^{-1}(u)]'du=\frac{1}{\lambda}\mu^{-1}(x)-\mu^{-1}(\mu_0)=\lambda\mu^{-1}(x).$$
		and since $\tilde{\mu}(t)=\Psi^{-1}(\lambda t)$ we have
		$$\lambda t=\Psi(\tilde{\mu}(t))=\lambda \mu^{-1}(\tilde{\mu}(t)),\;\;\Rightarrow \;\; \mu(t)=\tilde{\mu}(t),$$
		as we needed to show.
		
		
		
	\end{proof}

\end{document}
